\slajdid{03-s010}
\begin{frame}[t,label=03-s010]
\gusttekst
\fontsize{9}{9.5}\selectfont
\setlength{\parskip}{1pt}
Kod matrične strukture, treba nam dva selekciona signala na komponentama u slučaju dvodimenzionog dekodovanja, odnosno tru u slučaju trodimenzionog dekodovanja itd... Za veći broj izlaza matrična struktura će imati prednost zbog manjeg broja upotrebljenih komponenti, ali kao što smo rekli potrebne su specifičnije komponente. U tom smislu da vidimo kako bi izgledala realizacija dekodera \(6/2^6\) realizovana sa dekoderima 2/4.
\begin{columns}[c,onlytextwidth]
\begin{column}{.60\textwidth}\centering
\begingroup
\newcommand{\MDek}[3]{%
 \begin{scope}[shift={(#2,#3)}]
 \draw[DEvod] (0,0) rectangle(1.6,2.4);
 \node[align=center,inner sep=1pt] at (.8,1.90) {Dekoder\\2/4};
 \foreach \n/\y in {3/1.6,2/1.3,1/1.0,0/.7}{
  \node[anchor=east,inner sep=1pt] at (1.6,\y) {Y\n};
  \draw[DEvod] (1.6,\y)--(1.8,\y);
  \coordinate (#1-y\n) at (1.8,\y);
 }
 \foreach \n/\y in {1/1.0,2/.6}{
  \node[anchor=west,inner sep=1pt] at (0,\y) {CS\n};
  \draw[DEvod] (-.2,\y)--(0,\y);
  \coordinate (#1-cs\n) at (-.2,\y);
 }
 \foreach \n/\x in {1/.45,0/1.15}{
  \node[anchor=south,inner sep=1pt] at (\x,0) {S\n};
  \draw[DEvod] (\x,0)--(\x,-.2);
  \coordinate (#1-s\n) at (\x,-.2);
 }
 \end{scope}}
\begin{tikzpicture}[x=.86cm,y=.78cm,font=\fontsize{8}{9}\selectfont]
\MDek{v}{0}{2.8}
\node at (.27,4.2) {i};
\foreach \n/\y in {3/4.4,2/4.1,1/3.8,0/3.5}{
 \draw[DEvod] (v-y\n)--(3.75,\y);
 \node[anchor=south,inner sep=1pt] at (2.13,\y) {Y\n V};
}
\draw[DEvod] (1.0,0) rectangle (4.1,1.45);
\node[align=center,inner sep=1pt] at (1.70,.75) {Dekoder\\2/4};
\node at (3.4,.7) {j};
\foreach \n/\x in {3/2.65,2/3.0,1/3.35,0/3.7}{
 \draw[DEvod] (\x,1.45)--(\x,4.55);
 \node[rotate=90,anchor=east,inner sep=1pt] at (\x,1.45) {Y\n};
 \node[rotate=90,anchor=west,inner sep=1pt] at (\x,1.85) {Y\n K};
}
\foreach \n/\x in {1/2.55,2/3.0}{
 \node[rotate=90,anchor=west,inner sep=1pt] at (\x,0) {CS\n};
 \draw[DEvod] (\x,0)--(\x,-.12);
}
\node[rotate=90,anchor=south,inner sep=1pt] at (4.1,.40) {S1};
\node[rotate=90,anchor=south,inner sep=1pt] at (4.1,1.02) {S0};
\draw[DEvod] (4.1,.40)--(4.35,.40) node[right,inner sep=1pt] {S3};
\draw[DEvod] (4.1,1.02)--(4.35,1.02) node[right,inner sep=1pt] {S2};
\draw[DEvod] (v-cs1)--(-.35,3.8)--(-.35,-.22)--(2.55,-.22)--(2.55,-.12);
\draw[DEvod] (v-cs2)--(-.60,3.4)--(-.60,-.48)--(3.0,-.48)--(3.0,-.12);
\draw[DEvod] (-.35,-.22)--(-.8,-.22) node[left,inner sep=1pt] {CS1};
\draw[DEvod] (-.60,-.48)--(-.8,-.48) node[left,inner sep=1pt] {CS2};
\draw[DEvod] (v-s1)--++(0,-.25) node[below] {S5};
\draw[DEvod] (v-s0)--++(0,-.25) node[below] {S4};
\MDek{f}{4.8}{3}
\foreach \n in {3,2,1,0}
 \draw[DEvod] (f-y\n)--++(.08,0) node[right,inner sep=1pt] {Y(16i+4j+\n)};
\draw[DEvod] (f-cs1)--++(-.15,0) node[left,inner sep=1pt] {YiV};
\draw[DEvod] (f-cs2)--++(-.15,0) node[left,inner sep=1pt] {YjK};
\draw[DEvod] (f-s1)--++(0,-.25) node[below] {S1};
\draw[DEvod] (f-s0)--++(0,-.25) node[below] {S0};
\draw[DEvod,fill=white] (3.0,4.1) circle(.075);
\draw[DEstrelica] (3.0,4.18) .. controls (3.6,4.8) and (4.0,4.9) .. (4.5,4.9);
\end{tikzpicture}
\endgroup

\end{column}
\begin{column}{.37\textwidth}
U izlaznom nivou se nalazi 16 dekodera 2/4 (na slici je prikazan smo jedan) dok se jedan dekoder koristi za selekciju vrste (i) a drugi za selekciju kolone (j). Potreban broj komponenti nam je 16+1+1=18. Način povezivanja internih selekcionih signala je identičan kao I kod piramidalne strukture. Uočiti da bi nam kod piramidalne sturkture trebalo 16+4+1=21 komponenta.
\end{column}
\end{columns}

Ova razlika postaje još značajnija kao su u pitanju komponente sa još većim brojem izlaza. Takođe uočiti da će matrična struktura imati manje kašnjenje, pošto u ovom slučaju uvek su u pitanju dva nivoa, dok kod piramidalne strukture će broj nivoa zavisiti od veličine komponente. Ali za matričnu strukturu su potrebne komponenet sa više selekcionih signala. Normalno moguće je praviti i kombinacije ove dve strukture.

\end{frame}
